% Adenda v2 a Metodología: describe la campaña de CoppeliaSim que Resultados
% 6.6 reporta (mundo pareado de paso estrecho, E4), distinta de la escena de
% almacén AWS descrita en 04-methodology.tex. No edita ese archivo.
\subsection{CoppeliaSim: mundo pareado de paso estrecho (E4)}
\label{sec:method-coppelia-narrow}

Además del piloto de almacén de la Sección~\ref{sec:method_statistics_coppelia},
una segunda escena de CoppeliaSim reproduce el juego de acoplamiento de E4 en
un mundo de paso estrecho, para contrastar en un motor de simulación 3D
independiente la misma conclusión que el modelo cinemático de Python. La
escena real (no sintética) enfrenta cuatro AMR por método a través de un
cuello de botella entre dos obstáculos fijos, con origen, destino y semilla
(886001) idénticos para ambos métodos comparados: control directo con HOCBF y
el juego PD distribuido con HOCBF de E4. El protocolo registra la trayectoria
comandada y la reejecutada en el motor, y admite la campaña solo si supera
seis comprobaciones: escena guardada, cuadros reales capturados, al menos
treinta objetos en escena, error de reejecución por debajo de 1~mm, y
reproducción tanto del contraejemplo del método directo como del éxito seguro
del método propuesto. Como en el piloto de almacén, la evidencia se limita a
geometría y \emph{replay} cinemático; no sustituye contacto, dinámica de
ruedas ni validación de hardware.
